\documentclass[a4paper]{article}

% PACKAGES %%%%%%%%%%%%%%%%%%%%%%%%%%%%%%%%%%%%%%%%%%%%%%%%%%%%%%%%%%%%%%%%%%%%%
\usepackage[utf8]{inputenc}
\usepackage[english]{babel}
\usepackage{url}
\usepackage{hyperref}
\usepackage{charter}
\usepackage{float}
\usepackage{fancyhdr}
\usepackage{booktabs}
\usepackage{tabularx}
\usepackage{multirow}
\usepackage{enumitem}
\usepackage{xcolor}
\usepackage{graphicx}
\usepackage{siunitx}
\usepackage[hmarginratio=1:1, top=32mm, margin=0.8in, headsep=20mm]{geometry}

\newcommand{\TODO}{\textbf{\textcolor{red}{TODO}}}

\lhead{\includegraphics[scale=0.4]{CETC.png}}
\rhead{\includegraphics[scale=0.025]{Cornell-University-Logo.png}}

\pagestyle{fancy}
\fancyfoot{}
\fancyfoot[C]{\thepage}

\begin{document}

\thispagestyle{fancy}

\begin{center}
    \textbf{\large Orbit Odyssey -- Computer Vision Autonomous Challenge} \\[0.3em]
    {\Large Quotation}

    \vspace{1em}

    David Alejandro Fuquen Fl\'orez \\[0.3em]
    {\small \today}
\end{center}

\vspace{1em}

%===============================================================================
\section*{Summary}
%===============================================================================

This quotation covers the remaining development work for the Autonomous Computer Vision portion of the Orbit Odyssey Challenge, originally proposed by Rushil Sambangi (Cornell University). The scope includes three deliverables: a complete competition manual, adaptation of the graphical training interface, and refinement of the XRP robot control code.

The project will be delivered in three sequential phases over approximately three weeks.

%===============================================================================
\section*{Hourly Rate}
%===============================================================================

\begin{center}
\begin{tabular}{ll}
\toprule
\textbf{Hourly rate} & \$28.00 USD \\
\textbf{Estimated total hours} & 43 hours (base) \\
\textbf{Estimated total cost} & \$1,204.00 USD (base) \\
\textbf{Contingency range} & Up to 51 hours / \$1,428.00 USD \\
\bottomrule
\end{tabular}
\end{center}

\vspace{0.5em}
\noindent The contingency accounts for integration complexity with the existing CETC codebase, which has not yet been fully assessed. See Section ``Risk and Contingency'' for details.

%===============================================================================
\section*{Phase 1: Competition Manual}
%===============================================================================

\noindent\textbf{Timeline:} Week 1 \hfill \textbf{Estimated hours:} 14 \hfill \textbf{Cost:} \$392.00

\vspace{0.5em}

The manual will serve as the single source of truth for the competition. It will be inspired by and expand upon the original \href{https://docs.google.com/document/d/1N_mCX1EM3w5nDMbHrX3vcQlSqa2yJL38nCU0xxMPxmY/edit?tab=t.0}{Orbit Odyssey Manual}, with a strong emphasis on preventing cheating and closing loopholes.

\subsection*{Scope of Work}

\begin{enumerate}[leftmargin=*]
    \item \textbf{Playing Field Definition} (3--4 hrs)
    \begin{itemize}
        \item Field dimensions, boundaries, and markings
        \item Robot starting zone placement rules
        \item Object quantity, placement templates, and clearance constraints
        \item Object specifications (dimensions, materials, image faces)
    \end{itemize}

    \item \textbf{Scoring System and Scorecards} (3--4 hrs)
    \begin{itemize}
        \item Complete point system (field performance, mAP bonus, model-size multiplier)
        \item Penalty definitions and edge-case rulings
        \item Tiebreaker hierarchy
        \item Printable scorecard template for referees
        \item Integration of cryptographic receipt verification into scoring workflow
    \end{itemize}

    \item \textbf{Student Guide} (3--4 hrs)
    \begin{itemize}
        \item Step-by-step instructions for using the training GUI
        \item Explanation of hyperparameter choices and their trade-offs
        \item How to submit the trained model and receipt for deployment
        \item Rules of engagement (what teams can and cannot modify)
    \end{itemize}

    \item \textbf{Moderator Guide} (2--3 hrs)
    \begin{itemize}
        \item Dataset distribution procedure
        \item Field setup checklist (per-match randomization)
        \item Deployment pipeline: model upload, serial connection, run execution
        \item Receipt verification procedure (decryption and validation)
    \end{itemize}

    \item \textbf{Anti-Cheating Documentation} (1--2 hrs)
    \begin{itemize}
        \item How the signed receipt system works (for judges)
        \item What constitutes a valid vs.\ tampered receipt
        \item Disqualification criteria
    \end{itemize}
\end{enumerate}

\subsection*{Hour Breakdown}

\begin{center}
\begin{tabular}{lcc}
\toprule
\textbf{Task} & \textbf{Hours (est.)} & \textbf{Cost} \\
\midrule
Playing field definition & 3.5 & \$98.00 \\
Scoring system \& scorecards & 3.5 & \$98.00 \\
Student guide & 3.5 & \$98.00 \\
Moderator guide & 2.5 & \$70.00 \\
Anti-cheating documentation & 1.0 & \$28.00 \\
\midrule
\textbf{Phase 1 Total} & \textbf{14.0} & \textbf{\$392.00} \\
\bottomrule
\end{tabular}
\end{center}

%===============================================================================
\section*{Phase 2: Graphical User Interface}
%===============================================================================

\noindent\textbf{Timeline:} Weeks 2--3 \hfill \textbf{Estimated hours:} 18 \hfill \textbf{Cost:} \$504.00

\vspace{0.5em}

This phase adapts the existing training GUI (\texttt{gui\_train\_model.py}) to work with CETC's current codebase. It also implements a cryptographic anti-cheating system: a signed, encrypted receipt that records the exact hyperparameters and results from each training run. Only moderators possess the decryption key.

\subsection*{Scope of Work}

\begin{enumerate}[leftmargin=*]
    \item \textbf{CETC Code Integration} (5--12 hrs)
    \begin{itemize}
        \item Analyze CETC's existing training/data pipeline
        \item Adapt GUI to interface with CETC's code structure
        \item Ensure compatibility with CETC's dataset format and workflow
        \item Resolve any dependency or API conflicts
    \end{itemize}

    \item \textbf{Signed Receipt System} (4--6 hrs)
    \begin{itemize}
        \item Implement Fernet symmetric encryption with a moderator-held secret key
        \item Receipt captures: hyperparameters (epochs, learning rate, model size, train split), class names, final mAP per class, model file size, timestamp, and unique run ID
        \item Receipt exported as encrypted \texttt{.receipt} file alongside model weights
        \item Moderator-side decryption utility for verification
    \end{itemize}

    \item \textbf{mAP and Model-Size Extraction} (2--3 hrs)
    \begin{itemize}
        \item Programmatically extract mAP from training results CSV
        \item Compute and record model file size (bytes)
        \item Feed both values into the receipt and display in GUI summary
    \end{itemize}

    \item \textbf{UI Adjustments and Testing} (3--5 hrs)
    \begin{itemize}
        \item Update dataset selection flow for CETC's folder structure
        \item Add receipt export status indicator in GUI
        \item End-to-end testing: train $\rightarrow$ export $\rightarrow$ receipt generation $\rightarrow$ decryption verification
    \end{itemize}
\end{enumerate}

\subsection*{Hour Breakdown}

\begin{center}
\begin{tabular}{lcc}
\toprule
\textbf{Task} & \textbf{Hours (est.)} & \textbf{Cost} \\
\midrule
CETC code analysis \& integration & 8.0 & \$224.00 \\
Signed receipt system (Fernet) & 5.0 & \$140.00 \\
mAP / model-size extraction & 2.0 & \$56.00 \\
UI adjustments \& testing & 3.0 & \$84.00 \\
\midrule
\textbf{Phase 2 Total} & \textbf{18.0} & \textbf{\$504.00} \\
\bottomrule
\end{tabular}
\end{center}

\vspace{0.5em}
\noindent\textit{Note:} The CETC integration estimate carries the highest uncertainty. If their codebase requires significant refactoring to interface with the GUI, this task may require up to 12 hours (adding \$112 to Phase 2).

%===============================================================================
\section*{Phase 3: XRP Robot Code}
%===============================================================================

\noindent\textbf{Timeline:} Week 3 \hfill \textbf{Estimated hours:} 11 \hfill \textbf{Cost:} \$308.00

\vspace{0.5em}

Once the playing field and game rules are finalized (Phase 1), the XRP robot's MicroPython control code (\texttt{main.py}) will be adapted to operate within the defined arena constraints.

\subsection*{Scope of Work}

\begin{enumerate}[leftmargin=*]
    \item \textbf{State Machine Adaptation} (4--6 hrs)
    \begin{itemize}
        \item Align search/center/contact/avoid behavior with defined field dimensions
        \item Add timeout protection (prevent infinite search loops)
        \item Parameterize hardcoded values (distances, speeds, delays)
        \item Handle edge cases: object lost mid-approach, multiple objects in view
    \end{itemize}

    \item \textbf{Parameter Tuning and Testing} (3--5 hrs)
    \begin{itemize}
        \item Tune turn speeds, centering threshold, approach distances for physical field
        \item Test with actual objects (cans, wheels, rocket model)
        \item Validate rangefinder accuracy at competition distances
    \end{itemize}

    \item \textbf{Integration Testing} (2--4 hrs)
    \begin{itemize}
        \item Full pipeline test: GUI-trained model $\rightarrow$ \texttt{run\_model.py} $\rightarrow$ XRP
        \item Verify ACK flow control stability over full 120-second run
        \item Confirm scoring conditions are met (3~cm nudge for contact)
    \end{itemize}
\end{enumerate}

\subsection*{Hour Breakdown}

\begin{center}
\begin{tabular}{lcc}
\toprule
\textbf{Task} & \textbf{Hours (est.)} & \textbf{Cost} \\
\midrule
State machine adaptation & 5.0 & \$140.00 \\
Parameter tuning \& testing & 3.5 & \$98.00 \\
Integration testing (full pipeline) & 2.5 & \$70.00 \\
\midrule
\textbf{Phase 3 Total} & \textbf{11.0} & \textbf{\$308.00} \\
\bottomrule
\end{tabular}
\end{center}

%===============================================================================
\section*{Total Cost Summary}
%===============================================================================

\begin{center}
\begin{tabular}{lccc}
\toprule
\textbf{Phase} & \textbf{Hours} & \textbf{Cost} & \textbf{Timeline} \\
\midrule
1. Competition Manual & 14.0 & \$392.00 & Week 1 \\
2. GUI Adaptation & 18.0 & \$504.00 & Weeks 2--3 \\
3. XRP Robot Code & 11.0 & \$308.00 & Week 3 \\
\midrule
\textbf{Base Total} & \textbf{43.0} & \textbf{\$1,204.00} & \textbf{3 weeks} \\
\midrule
Contingency (CETC risk) & +8.0 & +\$224.00 & --- \\
\midrule
\textbf{Maximum Total} & \textbf{51.0} & \textbf{\$1,428.00} & \textbf{3 weeks} \\
\bottomrule
\end{tabular}
\end{center}

%===============================================================================
\section*{Weekly Milestones}
%===============================================================================

\begin{description}[leftmargin=!,labelwidth=5em]
    \item[Week 1] Complete competition manual. Deliver: finalized PDF with playing field specs, scoring rules, student/moderator guides, anti-cheat documentation, and printable scorecards.

    \item[Week 2] CETC integration and receipt system. Deliver: working GUI that produces encrypted receipts alongside trained models. Moderator decryption utility functional.

    \item[Week 3] XRP code finalization and full-pipeline integration testing. Deliver: updated \texttt{main.py}, tested end-to-end with physical hardware and competition objects.
\end{description}

%===============================================================================
\section*{Risk and Contingency}
%===============================================================================

\begin{enumerate}[leftmargin=*]
    \item \textbf{CETC Integration (HIGH risk):} The CETC codebase has not been fully assessed. If significant architectural differences exist between their pipeline and the current GUI, integration time could increase by 4--8 hours. Mitigation: early-week-2 assessment with go/no-go decision.

    \item \textbf{Hardware Availability (MEDIUM risk):} Phase 3 requires physical access to the XRP robot, webcam, and serial connection for testing. Delays in hardware access will compress the testing window.

    \item \textbf{Object Procurement (LOW risk):} Physical competition objects (cans, wheels, rocket model) must be available for Phase 3 testing. Testing can proceed with proxy objects if final versions are not yet ready.
\end{enumerate}

%===============================================================================
\section*{Assumptions}
%===============================================================================

\begin{itemize}
    \item Availability of 10--15 hours per week dedicated to this project.
    \item CETC codebase will be provided at the start of Phase 2.
    \item Hardware (XRP, webcam, serial cable) available for Phase 3 testing.
    \item Competition objects or suitable proxies available for integration testing.
    \item The competition uses Option 1 (open-field search), not Option 2 (line-following).
\end{itemize}

%===============================================================================
\section*{Payment Terms}
%===============================================================================

Payment may be structured per-phase or upon project completion, as agreed with the team. Hours will be tracked and reported weekly.

\end{document}
